\documentclass[10pt,a4paper]{article}

\usepackage[T1]{fontenc}
\usepackage[utf8]{inputenc}
\usepackage{lmodern}
\usepackage[margin=2.2cm]{geometry}
\usepackage{xcolor}
\usepackage{listings}
\usepackage{booktabs}
\usepackage{tabularx}
\usepackage{enumitem}
\usepackage{amsmath}
\usepackage{mathpartir}
\mprset{rightstyle=\normalfont\small}

% judgements
\newcommand{\eqj}[4]{\models #1 \sim #2 : #3 \Rightarrow #4}
\newcommand{\hj}[3]{\models #1 : #2 \Rightarrow #3}
\newcommand{\phj}[4]{\models #1 : #2 \Rightarrow #3\ [{=}\,#4]}
\newcommand{\Eq}[1]{\mathord{=}\{#1\}}
\newcommand{\cd}[1]{\mathtt{#1}}
\newcommand{\mods}{\mathrm{mod}}
\usepackage[colorlinks=true,linkcolor=blue!50!black,urlcolor=blue!50!black]{hyperref}

\setlist{nosep,leftmargin=*}
\setlength{\parskip}{0.5em}

% --------------------------------------------------------------------
% EasyCrypt listings
\lstdefinelanguage{easycrypt}{
  morekeywords=[1]{forall,exists,fun,glob,let,in,var,proc,if,then,else,elif,while,return,res,equiv,hoare,phoare,islossless},
  morekeywords=[2]{axiom,lemma,proof,qed,abort,end,import,export,include,local,declare,module,op,pred,require,theory,section,type,clone,with,Pr},
  morekeywords=[3]{simplify,congr,split,left,right,case,pose,have,elim,exlim,ecall,apply,rewrite,subst,progress,trivial,auto,move,transitivity,symmetry,seq,wp,sp,sim,skip,call,rcondt,rcondf,swap,rnd,bypr,byphoare,byequiv,conseq,inline,exists,first,last,do},
  morekeywords=[4]{exact,smt,by,done},
  morekeywords=[5]{admit},
  sensitive=true,
  morecomment=[n]{(*}{*)},
  morestring=[b]",
}

\lstset{
  language=easycrypt,
  upquote=true,
  columns=fullflexible,
  keepspaces=true,
  basicstyle=\footnotesize\ttfamily,
  keywordstyle=[1]{\bfseries\color{violet}},
  keywordstyle=[2]{\bfseries\color{olive!80!black}},
  keywordstyle=[3]{\bfseries\color{blue!70!black}},
  keywordstyle=[4]{\bfseries\color{teal}},
  keywordstyle=[5]{\bfseries\color{red}},
  commentstyle=\itshape\color{gray!80!black},
  frame=tb,
  numbers=left,
  numberstyle=\tiny\color{gray},
  numbersep=6pt,
  breaklines=true,
  literate={->}{{$\rightarrow$}}2 {=>}{{$\Rightarrow$}}2 {==>}{{$\Longrightarrow$}}3
           {/\\}{{$\wedge$}}2 {\\/}{{$\vee$}}2 {<=}{{$\le$}}2 {<>}{{$\neq$}}2,
}

\newcommand{\ec}{\lstinline[basicstyle=\small\ttfamily]}
\newcommand{\EC}{\textsc{EasyCrypt}}

\title{\EC{} tricks}
\author{}
\date{}

\begin{document}

\maketitle

\noindent
Every listing below is a complete file from this folder.
Each compiles on its own against the standard library, so any of them can be copied out, run, and changed.

\begin{lstlisting}[language={},numbers=none]
for f in *.ec; do easycrypt compile "$f" || echo "FAILED: $f"; done
\end{lstlisting}

\tableofcontents

% ====================================================================
\section{Rules of thumb}

\begin{itemize}
  \item \textbf{\ec|exlim| binds the value at the start of the remaining program.}
    That is the value in the current precondition.
    If the program changes it before the point where you need it, cut the program first with \ec|seq| or \ec|sp|.
    \ec|exlim e => v| is the same as \ec|exists* e; elim* => v|.
  \item \textbf{Calls are processed from the end.}
    When dropping one-sided calls, handle the last one first.
  \item \textbf{An equiv cannot relate a side to its own past.}
    Fix the old value with \ec|exlim| and add a hoare triple for that side with \ec|conseq E H1 H2|.
  \item \textbf{\ec|byequiv lemma| needs an exact match.}
    Write the spec you want in \ec|byequiv (: P ==> Q)| and then \ec|conseq lemma|.
    If the games are on the other sides, use \ec|symmetry| first.
  \item \textbf{An up-to-bad invariant must include the equality of the flag}, or nothing says the left flag is still down while the right one is.
  \item \textbf{A fact stated as $\Pr[\dots] = 1$ becomes a spec with \ec|bypr|.}
    For a phoare spec, \ec|bypr => &m <-; exact (h &m)| is enough.
    For a hoare spec, go through \ec|Pr[mu_not]| as in Section~\ref{sec:pr}.
\end{itemize}

% ====================================================================
\section{Reading the rules}\label{sec:reading}

Each section below states, as inference rules, the program-logic rule behind the tactics it shows.
The rules are simplified to the shape the tactic uses, and they are read bottom-up, because the tactic turns the goal below the line into the goals above it.

\begin{itemize}
  \item $\eqj{c_1}{c_2}{\Phi}{\Psi}$ is a pRHL judgement.
    From any two memories related by $\Phi$, the output distributions of $c_1$ and $c_2$ are related by $\Psi$, which means there is a coupling of them whose support lies in $\Psi$.
  \item $\Phi$ and $\Psi$ relate two memories.
    For a formula $P$ over one memory, $P\{1\}$ reads it in the left memory and $P\{2\}$ in the right one, and $\Eq{x}$ abbreviates $x\{1\} = x\{2\}$.
  \item $\hj{c}{P}{Q}$ is a Hoare judgement, saying that every terminating run of $c$ from a memory in $P$ ends in $Q$.
    It says nothing about termination.
  \item $\phj{c}{P}{Q}{p}$ is a probabilistic Hoare judgement, saying that from every memory in $P$ the probability that $c$ ends in $Q$ is exactly $p$.
    With $p = 1$ it gives both termination and the postcondition.
  \item $\mods(c)$ is the set of variables $c$ may write.
    When a rule quantifies over values for them, \EC{} shows the quantified names as \ec|result_L|, \ec|O_R|, \ec|S_R| and so on.
\end{itemize}


% ====================================================================
\section{\texorpdfstring{\texttt{conseq}}{conseq}}\label{sec:conseq}

\begin{itemize}
  \item \ec|conseq H| applies a spec lemma and leaves two side goals, the precondition implication and the postcondition implication.
  \item \ec|conseq (: _ ==> Q)| changes the postcondition to \ec|Q| and keeps the precondition.
    Use it to strengthen the postcondition to the equalities that \ec|sim| proves.
  \item \ec|conseq E H1 H2| conjoins a hoare triple for each side onto a two-sided equiv.
    An equiv relates the two sides to each other, never a side to its own initial value, so facts like ``the left counter went up by one'' have to come from a one-sided hoare triple.
  \item \ec|(_: true ==> true)| in place of \ec|H1| or \ec|H2| means that nothing is added on that side.
    It becomes a trivial goal.
  \item \ec|conseq ll H| turns losslessness and a hoare triple into a phoare spec with probability~$1$, which is what \ec|call{1}| needs.
  \item \ec|conseq H1 H2| on a hoare goal conjoins the postconditions of two hoare triples.
\end{itemize}

\begin{mathpar}
\inferrule*[right=\texttt{conseq}]
  {\Phi \Rightarrow \Phi' \\ \eqj{c_1}{c_2}{\Phi'}{\Psi'} \\ \Psi' \Rightarrow \Psi}
  {\eqj{c_1}{c_2}{\Phi}{\Psi}}

\inferrule*[right=\texttt{conseq E H1 H2}]
  {\eqj{c_1}{c_2}{\Phi'}{\Psi'} \\ \hj{c_1}{P_1}{Q_1} \\ \hj{c_2}{P_2}{Q_2} \\\\
   \Phi \Rightarrow \Phi' \wedge P_1\{1\} \wedge P_2\{2\} \\
   \Psi' \wedge Q_1\{1\} \wedge Q_2\{2\} \Rightarrow \Psi}
  {\eqj{c_1}{c_2}{\Phi}{\Psi}}

\inferrule*[right=\texttt{conseq ll H}]
  {\phj{c}{\mathit{true}}{\mathit{true}}{1} \\ \hj{c}{P}{Q}}
  {\phj{c}{P}{Q}{1}}
\end{mathpar}

\noindent
The second rule is sound because a Hoare postcondition holds in every final memory of its side, and so in every pair of memories the coupling can reach.
The last implication must hold in all such pairs, which is why \EC{} states it with the variables in $\mods(c_1)$ and $\mods(c_2)$ universally quantified.

\lstinputlisting[title={\texttt{01\_conseq.ec}}]{01_conseq.ec}

% ====================================================================
\section{\texorpdfstring{\texttt{call}}{call}}\label{sec:call}

\begin{itemize}
  \item \ec|call (: I)| handles a call to an abstract adversary with an invariant \ec|I|.
    It leaves one goal per oracle procedure, of the form \ec|={arg} /\ I ==> ={res} /\ I|.
    Directly on the procedure, the same thing is \ec|proc I|.
  \item A procedure declared with \ec|{}| in its module type may not call any oracle.
    \ec|call (: true)| is then enough, since it only touches the adversary's own state and every other fact in the precondition carries over.
  \item For an abstract oracle procedure, or an alias like \ec|proc get = O.get|, use \ec|proc*; call (: true)|.
  \item \ec|call (: I)| is only for abstract procedures.
    A concrete procedure needs a full spec, \ec|call (: P ==> Q)|.
  \item \ec|call{1} H| with \ec|H| a phoare spec with probability~$1$ runs the call on the left side only.
    Its postcondition is all you learn about what the call did.
  \item \ec|call H| with a hoare lemma \ec|H| works as in hoare logic.
\end{itemize}

\begin{mathpar}
\inferrule*[right=\texttt{call (: I)}]
  {\forall \cd{o}.\ \eqj{\cd{O}_1.\cd{o}}{\cd{O}_2.\cd{o}}{\Eq{\mathit{arg}} \wedge I}{\Eq{\mathit{res}} \wedge I}}
  {\eqj{\cd{A}(\cd{O}_1).\cd{f}}{\cd{A}(\cd{O}_2).\cd{f}}{\Eq{\mathit{arg}, \cd{glob\ A}} \wedge I}{\Eq{\mathit{res}, \cd{glob\ A}} \wedge I}}

\inferrule*[right=\texttt{call\{1\} H}]
  {\phj{\cd{f}}{P}{Q}{1} \\\\
   \Phi \Rightarrow P[e/\mathit{arg}]\{1\} \wedge
     \forall v\, \vec{w}.\ Q[v/\mathit{res}, \vec{w}/\mods(\cd{f})]\{1\} \Rightarrow \Psi[v/x, \vec{w}/\mods(\cd{f})]\{1\}}
  {\eqj{x \leftarrow \cd{f}(e)}{\mathtt{skip}}{\Phi}{\Psi}}
\end{mathpar}

\noindent
In the first rule $\cd{o}$ ranges over the oracle procedures $\cd{A}$ may call, $I$ must not mention $\cd{glob\ A}$, and $\cd{A}$ must not be able to write the variables of $I$.
A procedure declared with \ec|{}| may call no oracle, so the rule has no premises and $I = \mathit{true}$ suffices, while facts about variables outside $\mods = \cd{glob\ A}$ survive the call unchanged.

In the second rule the new values of everything $\cd{f}$ may modify are universally quantified, and $Q$ is all that is known about them.
A spec that says nothing about some variable the rest of the proof needs leaves an unprovable goal, with a fresh name such as \ec|S_R| where the old value used to be.

\lstinputlisting[title={\texttt{02\_call.ec}}]{02_call.ec}

% ====================================================================
\section{\texorpdfstring{\texttt{exlim}}{exlim} and \texorpdfstring{\texttt{seq}}{seq}}\label{sec:exlim}

\begin{itemize}
  \item \ec|exlim e => v| binds \ec|v| to the value of \ec|e| in the precondition of the current goal.
    That precondition describes the state at the start of the program that is left, not at the point where you want to use \ec|v|.
  \item The typical mistake is an \ec|exlim| at the top of a proof whose program changes the value before the call that needs it.
    The call's precondition then asks for an equality that no longer holds.
  \item The fix is to cut the program first, so that the call becomes the start of what is left.
    \ec|seq| cuts at any point, and \ec|sp| consumes leading assignments.
  \item In a relational goal, say which side, as in \ec|exlim x{1}, M.y{2} => a b|.
  \item An \ec|exlim| at the top is fine when nothing before the call changes the value, because then the value at the start and the value at the call coincide.
\end{itemize}

\begin{mathpar}
\inferrule*[right=\texttt{seq}]
  {\eqj{c_1}{c_2}{\Phi}{\Theta} \\ \eqj{c_1'}{c_2'}{\Theta}{\Psi}}
  {\eqj{c_1; c_1'}{c_2; c_2'}{\Phi}{\Psi}}

\inferrule*[right=\texttt{exlim e => v}]
  {\forall v.\ \eqj{c_1}{c_2}{\Phi \wedge e = v}{\Psi}}
  {\eqj{c_1}{c_2}{\Phi}{\Psi}}
\end{mathpar}

\noindent
In the second rule $v$ is fresh, and $e$ is evaluated in the memories $\Phi$ talks about, which are the initial memories of all of $c_1$ and $c_2$.
To name a value in the middle of a program, apply \ec|seq| first and then \ec|exlim| in its second premise, where the middle of the program has become the start of $c_1'$ and $c_2'$.

\lstinputlisting[title={\texttt{03\_exlim\_seq.ec}}]{03_exlim_seq.ec}

% ====================================================================
\section{Probabilities}\label{sec:pr}

\begin{itemize}
  \item \ec|rewrite Pr[lemma]| applies a distribution lemma to a probability term.
    The useful ones are \ec|mu_le1|, \ec|mu_sub|, \ec|mu_not|, \ec|mu_or| and \ec|mu_split E|.
  \item $\Pr[E] = 0$ follows from a hoare triple on $\neg E$, by \ec|byphoare|, then \ec|hoare|.
  \item A phoare spec with probability~$1$ follows from a fact of the form $\Pr[\dots] = 1$ that holds at every memory, by \ec|bypr|.
  \item A hoare spec follows from the same fact.
    \ec|bypr| turns the goal into $\Pr[\neg Q] = 0$, and $\Pr[\neg Q] = \Pr[\mathit{true}] - \Pr[Q]$ with $1 = \Pr[Q] \le \Pr[\mathit{true}] \le 1$.
    No losslessness is needed.
  \item Inside an event written in a proof, \ec|X{m}| for the memory bound by \ec|bypr| is read as the current \ec|X|.
    Rather than retyping the event, rewrite with the hypothesis that already contains it, as in \ec|rewrite -h1 Pr[mu_sub]|.
  \item \ec|byequiv lemma| needs the lemma's pre- and postcondition to match exactly.
    State the spec in \ec|byequiv|, discharge it with \ec|conseq|, and use \ec|symmetry| when the games sit on the other sides.
  \item \ec|smt| treats probability terms as opaque reals, so arithmetic over several of them works as long as the terms are syntactically identical.
\end{itemize}

\begin{mathpar}
\inferrule*[right=\texttt{byequiv}]
  {\eqj{c_1}{c_2}{\Phi}{\Psi} \\ \Phi(m_1, m_2) \\ \Psi \Rightarrow (E_1\{1\} \Rightarrow E_2\{2\})}
  {\Pr[c_1, m_1 : E_1] \le \Pr[c_2, m_2 : E_2]}

\inferrule*[right=\texttt{symmetry}]
  {\eqj{c_2}{c_1}{\Phi^{-1}}{\Psi^{-1}}}
  {\eqj{c_1}{c_2}{\Phi}{\Psi}}

\inferrule*[right=\texttt{bypr} (phoare)]
  {\forall m.\ P(m) \Rightarrow \Pr[c, m : Q] = p}
  {\phj{c}{P}{Q}{p}}

\inferrule*[right=\texttt{bypr} (hoare)]
  {\forall m.\ P(m) \Rightarrow \Pr[c, m : \neg Q] = 0}
  {\hj{c}{P}{Q}}

\inferrule*[right=\texttt{byphoare; hoare}]
  {\hj{c}{P}{\neg E} \\ P(m)}
  {\Pr[c, m : E] = 0}
\end{mathpar}

\noindent
With $\Leftrightarrow$ in place of $\Rightarrow$ in the last premise, \ec|byequiv| gives equality.
$\Phi^{-1}$ is $\Phi$ with $\{1\}$ and $\{2\}$ exchanged.

\lstinputlisting[title={\texttt{04\_pr.ec}}]{04_pr.ec}

% ====================================================================
\section{Up to bad}\label{sec:upto}

The call rule takes three components, \ec|call (: B, I, J)|.

\begin{center}
\begin{tabular}{@{}lp{.78\textwidth}@{}}
  \toprule
  \ec|B| & the bad event, read on the right memory and written without \ec|{2}| \\
  \ec|I| & the relational invariant while $\neg$\ec|B|, which must include the equality of the flag \\
  \ec|J| & what still holds once \ec|B| is set, preserved by each side alone (optional) \\
  \bottomrule
\end{tabular}
\end{center}

\noindent
It generates a goal saying that the adversary terminates when its oracles do, and then three goals for every oracle procedure \ec|o|.

\begin{enumerate}
  \item \ec|equiv [o ~ o : !B{2} /\ ={arg} /\ I ==> if B{2} then J else ={res} /\ I]|.
  \item \ec|forall &2, B{2} => phoare [o{1} : J ==> J] = 1%r|, for the left side after bad.
  \item \ec|forall &1, phoare [o{2} : B /\ J ==> B /\ J] = 1%r|, which says the right side keeps bad.
\end{enumerate}

\noindent
The code before the call is left with the postcondition \ec|if B{2} then J else (={res, glob A} /\ I)|.

\begin{itemize}
  \item Choose \ec|J| so that it implies bad stays set.
    If some oracle could unset bad, for instance by overwriting the record that bad is computed from, put into \ec|J| whatever switches that oracle off.
  \item If you only need \ec|!B{2} => ={res}| and not the equality of the flags afterwards, drop \ec|J|.
    The goals after bad then only ask for termination.
  \item From the equivalence, one \ec|byequiv| gives $\Pr_1[E] \le \Pr_2[E \vee \mathit{bad}] \le \Pr_2[E] + \Pr_2[\mathit{bad}]$.
  \item If $\Pr_2[\mathit{bad}] = 0$, which a hoare triple can show, a second \ec|byequiv| in the other direction gives $\Pr_2[E \wedge \neg\mathit{bad}] \le \Pr_1[E]$.
    Splitting $\Pr_2[E]$ on bad then yields equality.
\end{itemize}

\noindent
Writing $\mathcal{I}(X)$ for $\mathbf{if}\ B\{2\}\ \mathbf{then}\ J\ \mathbf{else}\ (\Eq{X} \wedge I)$, the rule reads as follows.

\begin{mathpar}
\inferrule*[right=\texttt{call (: B, I, J)}]
  {\text{$\cd{A}$ is lossless when its oracles are} \\\\
   \forall \cd{o}.\ \eqj{\cd{O}_1.\cd{o}}{\cd{O}_2.\cd{o}}{\neg B\{2\} \wedge \Eq{\mathit{arg}} \wedge I}{\mathcal{I}(\mathit{res})} \\\\
   \forall \cd{o}.\ \forall m_2.\ B(m_2) \Rightarrow \big(\phj{\cd{O}_1.\cd{o}}{J}{J}{1}\big) \\
   \forall \cd{o}.\ \forall m_1.\ \phj{\cd{O}_2.\cd{o}}{B \wedge J}{B \wedge J}{1}}
  {\eqj{\cd{A}(\cd{O}_1).\cd{f}}{\cd{A}(\cd{O}_2).\cd{f}}{\mathcal{I}(\mathit{arg}, \cd{glob\ A})}{\mathcal{I}(\mathit{res}, \cd{glob\ A})}}

\inferrule*[right=up to bad]
  {\eqj{c_1}{c_2}{\Phi}{\Eq{B} \wedge (\neg B\{2\} \Rightarrow \Eq{\mathit{res}})} \\ \Phi(m_1, m_2)}
  {\Pr[c_1, m_1 : E] \le \Pr[c_2, m_2 : E] + \Pr[c_2, m_2 : B]}
\end{mathpar}

\noindent
In the third and fourth premises $J$ is a relation, and the memory of the side that does not run is the fixed $m_2$ or $m_1$.
The second rule is \ec|byequiv| followed by $\Pr[E \vee B] \le \Pr[E] + \Pr[B]$, where $E$ is an event on the result.

\lstinputlisting[title={\texttt{05\_upto\_bad.ec}}]{05_upto_bad.ec}

% ====================================================================
\section{Dropping one-sided calls}\label{sec:onesided}

A rewinding reduction runs the adversary once more than the game it is compared with, and that extra run is dropped on one side only.

\begin{enumerate}
  \item For every call you drop, get a phoare spec with probability~$1$ that says what it does to the state you still care about.
    It comes from losslessness and a hoare fact through \ec|conseq ll H|, or from an axiom of the form $\Pr[\dots] = 1$ through \ec|bypr|.
    If the axiom is existential, as a rewinding axiom usually is, these specs must be \ec|have|s inside the proof, after \ec|elim|.
  \item Fix, with \ec|exlim|, the values those specs need, at a point where they are right.
  \item Run \ec|call{1}| with each spec, from the last dropped call backwards.
  \item Close with \ec|auto| and \ec|smt|.
\end{enumerate}

\noindent
The state-reading procedures are declared with \ec|{}|, so they cannot call the oracle.
They stand for the reduction taking a snapshot of the adversary, not for something the adversary does.

\noindent
Each dropped call is one application of \ec|call{1} H| from Section~\ref{sec:call}.
For a discarded run between a save and a restore, the derivation is the following chain, read from the bottom.

\begin{mathpar}
\inferrule*[right=\texttt{call\{1\} (phs g)}]
 {\inferrule*[right=\texttt{call\{1\} (phr k)}]
   {\inferrule*[right=\texttt{call\{1\} (phl g)}]
     {\Phi \wedge \cd{glob\ A}\{1\} = g \Rightarrow \Psi}
     {\eqj{\cd{load}(s)}{\mathtt{skip}}{s\{1\} = f(g) \wedge \Phi}{\Psi}}}
   {\eqj{\cd{run}(); \cd{load}(s)}{\mathtt{skip}}{s\{1\} = f(g) \wedge \Phi}{\Psi}}}
 {\eqj{s \leftarrow \cd{save}(); \cd{run}(); \cd{load}(s)}{\mathtt{skip}}{\cd{glob\ A}\{1\} = g \wedge \Phi}{\Psi}}
\end{mathpar}

\noindent
Here $\Phi$ mentions only state that $\cd{run}$ leaves alone, which is what the spec \ec|phr| provides, and $\cd{load}(f(g))$ restores $\cd{glob\ A}$ to $g$ whatever $\cd{run}$ did to it.

\lstinputlisting[title={\texttt{06\_one\_sided.ec}}]{06_one_sided.ec}

% ====================================================================
\section{Pitfalls}\label{sec:pitfalls}

\begin{itemize}
  \item \ec|forall| extends as far to the right as it can.
    In an invariant write \ec|/\ (forall y, P y) /\ Q| with the parentheses, or \ec|Q| ends up under the binder.
  \item \ec|split| and other tactic names cannot be used as lemma or hypothesis names.
  \item Program expressions cannot mention \ec|glob M|.
    \EC{} reports that expressions cannot contain a glob statement.
  \item Inside a section, a \ec|declare axiom| cannot use a \ec|local| predicate, and a non-local operator cannot mention the \ec|glob| of a declared module.
    Spell the predicate out in the axiom.
  \item \ec|=> //| may close more goals than expected, and the next bullet then fails because all goals are closed.
  \item \ec|ecall{1} (lemma e{1})| crashed with \ec|InvalidGoalShape| in the version used here.
    \ec|seq|, \ec|exlim| and \ec|call{1}| do the same job.
  \item The order of the side goals from \ec|conseq| is not always the order of its arguments.
    Look at them before writing bullets.
\end{itemize}

% ====================================================================
\section{Seeing the goals from the command line}\label{sec:llm}

\ec|easycrypt llm| loads a file up to a line and prints the goals there.

\begin{lstlisting}[language={},numbers=none]
easycrypt llm -I <dirs> -eval 'LOAD "file.ec" 42
QUIET ON
proc.
GOALS ALL'
\end{lstlisting}

\begin{itemize}
  \item \texttt{LOAD} checks everything up to and including the given line, then waits for more input.
    If that line is \ec|proof. admit. qed.|, split it first, or the load runs past the proof.
  \item Every further line is one tactic, one sentence per line.
  \item \texttt{GOALS} prints the focused goal, \texttt{GOALS ALL} prints all of them, and \texttt{TREE ALL} prints the goal tree.
    An error prints \texttt{ERROR} together with the offending \texttt{source:} line.
\end{itemize}

\end{document}
